\section{Evaluation}
\label{sec:evaluation}

We evaluate \sys's efficiency in tracking objects while maintaining tracking fidelity, comparing against prior systems on seven datasets.

\begin{figure*}[ht]
    \centering
    \includegraphics[width=\textwidth]{figures/generated/revision_hota_summary.pdf}
    \caption{End-to-end throughput-accuracy results across seven datasets.
    \panellabel{Top}
    Tracking accuracy (HOTA) versus throughput for \sys{}, the reference pipeline, prior systems, and ablations.
    \panellabel{Middle}
    Speedup of \sys{} over each prior system at matched HOTA accuracy levels.
    \panellabel{Bottom}
    HOTA gain of \sys{} over each prior system at matched throughput levels.
    % \revpending{LEAP uses the higher-HOTA valid reconstruction per dataset from \autoref{tab:leap-reconstruction}.}
    % \ck{split the ablation if we have space.}
    }
    \Description{}
    \label{fig:e2e-results}
\end{figure*}

% \paragraph
\textbf{\em Accuracy Metrics.}
Prior systems evaluate accuracy using predefined queries~\cite{kang2017noscope,bastani2022otif,xu2024leap},
which can restrict evaluation to anticipated situations and fail to capture more complex object behaviors.
We instead adopt HOTA (\autoref{sec:track-materialization}) as a query-agnostic measure of track reconstruction quality.
HOTA
% evaluates reconstructed tracks by matching predicted and ground-truth detections at multiple localization thresholds.
% It
balances detection accuracy with trajectory association accuracy,
so missing detections,
fragmented tracks,
and merged identities all reduce the score.

% \paragraph
\textbf{\em Datasets.}
We evaluate \sys{} on seven video datasets (CalDoT1/2, B3D1/2/3/4, and Amsterdam) that vary in scene density, camera viewpoint, and object motion patterns,
allowing us to test performance across a diverse set of real-world conditions.
CalDoT1/2 (720 $\times$ 480 pixels) were introduced by OTIF~\cite{bastani2022otif},
while Amsterdam (1280 $\times$ 720 pixels) was introduced by NoScope~\cite{kang2017noscope}.
We add the B3D (1080 $\times$ 720 pixels) datasets
(B3D1/2/3/4)
to expand the evaluation beyond existing benchmarks.
All datasets consist of approximately one-minute clips captured at 15 FPS;
CalDoT1/2,
and Amsterdam each include 180 clips,
while each B3D dataset includes 18 clips.

CalDoT1/2~\cite{bastani2022otif} consist of highway traffic footage from California Department of Transportation cameras;
Amsterdam~\cite{kang2017noscope} is captured from an urban street-side camera.
These datasets exhibit structured motion in relatively constrained settings.
The B3D datasets are drawn from the Berkeley DeepDrive Drone Dataset~\cite{wu2025b3d},
consisting of aerial footage of unsignalized intersections and roundabouts with less predictable motion and higher interaction density.
We mask non-road and parking areas in the B3D videos.
By including B3D, we extend beyond the structured highway and low-density roadway settings of CalDoT and Amsterdam.
% \alvin{what does robustness mean ad how is this more dynamic than the other datasets?}

For CalDoT1/2 and Amsterdam,
we partition the data into training,
validation,
and testing splits,
each consisting of 60 one-minute clips.
For the B3D datasets,
because of the amount of footage available,
we use training,
validation,
and testing splits of 8,
5,
and 5 approximately one-minute clips,
respectively.

% \ck{we can mention that caldots come from OTIF, might make our evaluation results stronger} \ck{amsterdam, if I remember correctly, is not from OTIF. So, amsterdam videos can represent independent videos that are not from either of the system (us / otif).}

% \paragraph
\textbf{\em Ground Truth.}
We obtain ground truth detection and tracking annotations for the test splits used in our end-to-end evaluation (\autoref{sec:e2e-results}).
For CalDoT1 and CalDoT2,
we use the dense per-frame object detection labels provided by the OTIF dataset.
For Amsterdam,
we start from the same OTIF dense labels,
filter to vehicle classes only,
and apply minor manual corrections to remove erroneous annotations.
For B3D1/2/3/4,
we generate dense detection labels using
\revadd{
Faster R-CNN~\cite{ren2016fasterrcnn} trained with training labels from the B3D paper.
The labels are then verified and corrected by the authors.}
% to ensure high-quality ground truth for evaluation.
% the RetinaNet~\cite{lin2017focal} detector from the B3D paper with overlapping tiled inference~\cite{unel2019tiling,akyon2022sahi}.
% Specifically,
% we extract four fixed crops at the corners of each frame,
% each spanning two-thirds of the frame width and height,
% and merge the resulting detections via non-maximum suppression (NMS).

Ground truth tracking annotations for all datasets are generated using OC-SORT~\cite{cao2023ocsort},
a multi-object tracker that supersedes BYTETrack~\cite{zhang2022bytetrack} and SORT~\cite{bewley2016sort}.
In our experiments, the user-provided tracker is BYTETrack,
which also serves as the reference while generating the Pareto frontier.
% \sys{}'s built-in SORT is the second tracker option considered during Pareto frontier generation.
OC-SORT is used only for ground truth generation,
ensuring that evaluation results are not biased by shared tracking components.
Ground-truth annotations are used only for evaluation;
\sys does not access them while learning the artifacts or generating the Pareto frontier.

% \paragraph
\textbf{\em Prior Systems.}
We compare \sys{} against two prior approaches: OTIF~\cite{bastani2022otif} and LEAP~\cite{xu2024leap}.
OTIF utilizes a lightweight proxy model to identify irrelevant regions within frames for more efficient batching, along with uniform frame sampling and a specialized tracker to further reduce calls to the object detector.

LEAP uses predictive sampling to prune frames that are unlikely to contain relevant information.
For both methods,
we use the authors' publicly available implementations.

For OTIF,
we use the training split for its recurrent rate tracker,
since the original ``tracker'' split is unavailable across all datasets.
For LEAP,
the distilled detector weights are not publicly available,
so we use the original non-distilled detector to ensure reproducibility.

% In particular, while the authors describe training on multiple datasets (e.g., PASCAL VOC, Stanford Car Dataset, Object365, and VisDrone), they retain only samples with “vehicle-related” labels without clearly defining this subset. 

For a fair comparison across systems,
we standardize the object detector used in all experiments.
Specifically,
we use YOLOv5~\cite{khanam2024yolov5} for the CalDoT1, CalDoT2, and Amsterdam datasets,
and RetinaNet~\cite{lin2017focal} for the B3D datasets,
i.e., the detector released with B3D.
This ensures that any throughput and accuracy differences arise from the system rather than underlying detection models.
% \revpending{For CalDoT1/2 and Amsterdam,
% we rerun LEAP's preparation and extraction with the same dataset-specific YOLOv5 checkpoint as the reference pipeline,
% regenerating its pseudo-labels,
% re-identification model,
% and motion patterns.}
% \revpending{The reference pipeline,
% \sys{},
% and the recorded OTIF configurations use the dataset-specific detector weights.
% LEAP's CalDoT/Amsterdam runs use its non-distilled YOLOv5x weights;
% its B3D runs use the released RetinaNet.
% This substitution limits attribution of LEAP's accuracy differences to sampling alone.}
For \sys{},
the user-provided tracker is BYTETrack and the built-in default is SORT,
so Pareto frontier generation searches over the two tracker options $\{\text{BYTETrack}, \text{SORT}\}$.
When generating the Pareto frontier (\autoref{sec:pareto-frontier}),
BYTETrack at full frame rate serves as the user-provided tracker in the
% full-frame,
% every-frame
reference pipeline~(\autoref{sec:tracking-by-detection}), and 
the Pareto-optimal operating points may draw from either tracker option.

% \paragraph
\textbf{\em Detector Training.}
We train the YOLOv5 models on the training splits using the same curated labels described above:
the dense detection labels from the OTIF dataset for CalDoT1 and CalDoT2,
and the vehicle-filtered,
manually corrected OTIF labels for Amsterdam.

% \paragraph
\textbf{\em Experiment Setup.}
All evaluated systems consist of two stages:
a preparation stage to learn from the training split and generate their configurations using the validation split,
and a track extraction stage to track objects in the test split.
% \revpending{We separate training-data preparation from test-split extraction.
% \sys{} and OTIF select configurations using validation;
% LEAP uses one fixed configuration.}
We report throughput on track extraction only.
\revadd{For \sys{},
this includes relevance classification,
polyomino grouping,
ILP construction and solving,
packing,
detection,
unpacking,
tracking,
and interpolation.
Preparation is accounted for separately in \autoref{sec:preparation-cost}.}
% and exclude preparation.
Each data point in the results corresponds to a distinct system configuration.
For \sys{} and OTIF,
configurations are the Pareto-optimal operating points selected on the validation set (\autoref{sec:pareto-frontier}).
% For OTIF,
% configurations are produced by its joint parameter tuning procedure~\cite{bastani2022otif}.
LEAP, having a single fixed configuration, contributes one point.
For \sys{} and OTIF,
we report only the Pareto-optimal points from the test split results.
% \alvin{LEAP doesn't have pareto points?}

% \paragraph
\textbf{\em Hardware.}
All experiments are conducted on a server equipped with dual Intel Xeon Gold 6248 CPUs (2.50\,GHz),
504\,GB of RAM,
and an NVIDIA Quadro RTX 6000 GPU (24\,GB VRAM).
% All systems are evaluated using a single GPU.
Reported runtimes measure computation time only,
excluding disk I/O.

% \paragraph
\textbf{\em Track Interpolation.}
Tracking outputs from all systems are linearly interpolated to produce detection points on every frame for every track.
This step ensures a fair comparison under the HOTA metric,
which penalizes missing per-frame detections within tracks.
Systems that employ frame skipping inherently produce tracks with detection gaps;
without interpolation,
these gaps would exaggerate the HOTA penalty even when the underlying motion is faithfully captured,
for example,
a vehicle traveling at constant speed along a straight path.
By applying interpolation uniformly across all systems,
we isolate tracking quality from the detection gaps that frame-sampling strategies introduce.
% \revpending{For \sys{} and OTIF,
% we linearly interpolate between consecutive observations of the same track,
% without extrapolating beyond its observed endpoints.
% This fills sampling gaps for straight,
% constant-speed motion but cannot recover arbitrary motion changes.
% For LEAP,
% we separately evaluate raw boxes,
% linear interpolation,
% and assigned-pattern reconstruction (\autoref{sec:leap-queries}).
% HOTA then measures both per-frame detection and association quality.}

\subsection{End-to-end Results}
\label{sec:e2e-results}

\autoref{fig:e2e-results} summarizes the throughput-accuracy tradeoff across all seven datasets.
% \sys{}'s Pareto frontier is produced by a configuration search that treats the tracking algorithm as one of its dimensions,
% so the frontier for each dataset may contain operating points drawn from both SORT and BYTETrack,
% depending on which algorithm yields better throughput or accuracy in that regime (\autoref{sec:candidate-algorithms}).
\sys{} consistently dominates the Pareto frontier,
achieving higher HOTA scores than both OTIF and LEAP at the same throughput,
and higher throughput at the same accuracy level.
At matched throughput,
we compare each prior work's Pareto-optimal configuration against the most accurate \sys{} configuration that is at least as fast.
\sys{} achieves higher HOTA in every such comparison across all seven datasets
(\compareMatchedTputReachableCount{} of \compareMatchedTputPriorConfigCount{} prior works' configurations).
Of the remaining \compareMatchedTputUnreachableCount{} prior configurations that exceed \sys{}'s highest throughput,
\compareMatchedTputUnreachableZeroHotaCount{} produce zero HOTA;
the last achieves a \compareMatchedTputUnreachableRestSpeedup{}$\times$ speedup over \sys{}'s fastest configuration
while dropping HOTA by \compareMatchedTputUnreachableRestHotaDrop{}
on the \compareMatchedTputUnreachableRestDataset{} dataset.
The median HOTA gain across these comparisons is \compareMatchedTputMedianHotaGain{},
per-dataset maximum gains range from \compareMatchedTputPerDatasetMaxGainMin{} to \compareMaxHotaImprovement{},
and the largest gain occurs at \compareMaxHotaImprovementFps{} FPS on the \compareMaxHotaImprovementDataset{} dataset.

In realistic deployment scenarios where tracking accuracy is critical,
we measure the maximum speedup each system achieves over the full-frame,
every-frame reference pipeline at bounded HOTA accuracy loss.
% \todo{Add speedup table.}
\sys{} achieves a speedup while staying within a 5\% HOTA drop from the reference pipeline on all
% \comparePolytrisMeetFivePct{}/7
\revisionDatasetCount{} datasets,
while prior systems meet this constraint on only \comparePriorMeetFivePct{}/\revisionDatasetCount{} datasets.
% ,
% failing on \comparePriorFailDatasetsFivePct{} datasets.
To avoid overclaiming \sys{}'s improvement over prior systems,
we restrict prior systems to Pareto-optimal configurations and compare against a prior configuration whose HOTA is no higher than \sys{}'s selected configuration.
% \revpending{To avoid overclaiming,
% we first choose the fastest \sys{} configuration within the loss bound.
% If the fastest Pareto-optimal prior configuration meeting the bound has HOTA no higher than this \sys{} point,
% we use it.
% Otherwise,
% we use the most accurate prior Pareto point with HOTA no higher than \sys{}'s,
% breaking accuracy ties by throughput.
% The bound therefore constrains \sys{};
% the compared prior point can fall below it.}
Under this 5\% constraint,
% \revpending{With \sys{} constrained to at most 5\% HOTA loss,}
\sys{} achieves \compareAccuracyMatchedSpeedupMinFivePct{}$\times$ to \compareAccuracyMatchedSpeedupMaxFivePct{}$\times$ speedup over prior systems,
and \compareNaiveSpeedupMinFivePct{}$\times$ to \compareNaiveSpeedupMaxFivePct{}$\times$ speedup over the reference pipeline.
At a 10\% HOTA drop,
% \revpending{Applying the same selection rule with a 10\% loss bound for \sys{},}
\sys{} achieves \compareAccuracyMatchedSpeedupMinTenPct{}$\times$ to \compareAccuracyMatchedSpeedupMaxTenPct{}$\times$ speedup over prior systems,
and \compareNaiveSpeedupMinTenPct{}$\times$ to \compareNaiveSpeedupMaxTenPct{}$\times$ speedup over the reference pipeline.
Under both constraints,
% \revpending{In both comparisons,}
\sys{} consistently achieves higher speedup than all other systems on every dataset.

% \paragraph
\textbf{\em Analysis.} LEAP is designed to optimize for object-existence queries,
such as counting or detecting the time period during which an object is present,
rather than reconstructing full per-frame object tracks.
Accordingly,
LEAP outputs a time interval and \revadd{an assigned} motion pattern for each detected object rather than a per-frame trajectory.
% \revdelblock{To evaluate LEAP under the HOTA metric,
% we assign the closest reference track to each detected object.
% However,
% the resulting trajectories exhibit substantial misalignment with ground truth tracks,
% leading to low accuracy scores across all datasets.}
\revadd{For LEAP's HOTA evaluation,
we consider two approaches:
% \revdel{(i) raw boxes on sampled frames,}
(i) linear interpolation between consecutive boxes of the same object on sampled frames,
and (ii) trajectories reconstructed from LEAP's assigned motion patterns and predicted presence intervals.
The second approach uses LEAP's own pattern assignments without matching objects to evaluation reference tracks.
% \revdel{We report all three HOTA scores in Table~2 and use}
We use the highest-HOTA approach for each dataset in comparisons with \sys{} and OTIF,
giving LEAP its most favorable result among the two approaches.}

\revdelblock{Dense B3D scenes contain varied motion and frequent object entries and exits,
which can limit the coverage of a fixed motion-pattern library and leave few observations of each object.
For example,
LEAP observes \revisionLeapBThreeDThreeObjectsOnce{} of its \revisionLeapBThreeDThreeObjects{} predicted objects on B3D3 only once,
leaving no pair of boxes to interpolate.
An assigned motion pattern can also differ from an object's actual path,
so an approximately correct presence interval need not yield well-aligned per-frame boxes or high HOTA.}
\revpending{LEAP's low HOTA reflects masking,
association,
and reconstruction losses.
Its configured masks completely obscure \revisionLeapCalDoTOneMaskedBoxesPct{}\% and \revisionLeapCalDoTTwoMaskedBoxesPct{}\% of evaluated ground-truth boxes on CalDoT1 and CalDoT2,
respectively.
On CalDoT1,
only \revisionLeapCalDoTOneMovingContinuityPct{}\% of consecutive ground-truth-matched observations below LEAP's stationary-overlap threshold retain an identity.
On CalDoT2,
\revisionLeapCalDoTTwoEarlyExitPct{}\% of uncensored tracks lasting at least one second and having a following sample exit between their first two recorded detection frames.
Even with true boxes and identities at LEAP's sample times,
linear interpolation reaches IoU $\geq0.5$ on only \revisionLeapCalDoTTwoOracleBoxMatchPct{}\% of bracketed CalDoT2 boxes;
straight roads do not imply uniform image-space motion.
Assigned-pattern reconstruction also aligns with only \revisionLeapCalDoTTwoPatternAnchorMatchPct{}\% of LEAP's observed boxes at this threshold.}

\revpending{B3D's unsignalized intersections and traffic circles also expose pattern-coverage and association limitations.
On B3D3,
\revisionLeapBThreeDThreeEligiblePatterns{}/\revisionLeapBThreeDThreePatterns{} learned patterns pass the configured endpoint-region filter,
and only \revisionLeapBThreeDThreeContinuityPct{}\% of consecutive ground-truth-matched observations retain an identity.
Thus repeated observations can still yield separate one-observation identities.
Amsterdam's higher reported HOTA uses linear reconstruction;
its limited movement between many sampled endpoints makes interpolation more accurate,
with the same true-box diagnostic reaching \revisionLeapAmsterdamOracleBoxMatchPct{}\% at IoU $\geq0.5$.}

% \input{content/revision/native_queries}
% \input{content/revision/leap_reconstruction}

On the other hand, OTIF achieves reasonable accuracy,
but its speedup derives primarily from uniform frame skipping,
with its segmentation proxy model contributing comparatively little to throughput gains,
as reported by Bastani et al.~\cite{bastani2022otif}.
This reliance on frame skipping produces coarse-grained tracks that degrade accuracy at higher throughput operating points,
consistent with the observation that uniform temporal sampling cannot accommodate spatially varying sampling requirements (\autoref{sec:motivation}, Obs.~2).
The limited benefit of OTIF's proxy is visible directly in the B3D4 Pareto front.
Its most accurate point disables the segmentation proxy entirely,
while the next point enables it and exhibits a sharp accuracy drop with little speedup gain.
This points to a limitation that motivates \sys{}'s design:
OTIF's rectangular crops are a poor fit for the irregular shape of real activity regions (\autoref{sec:motivation}, Obs.~3),
forcing the proxy to either over-include background (yielding little speedup) or aggressively discard tiles with detections (yielding the observed accuracy drop).
% Second, the proxy considers only the pixel content of the current frame,
% ignoring the motion and positional cues that \sys{}'s relevance classifier combines;
% without these cues, the proxy cannot reliably distinguish moving objects from background,
% precisely the failure mode that \autoref{sec:relevance-classification} aims to addresses.

\input{content/revision/preparation}

\subsection{Ablation Studies}

We now systematically break down the contribution of the operators in \sys{}'s tracking pipeline.
We examine the throughput improvements achieved by sequentially enabling our core spatial and temporal optimization operators,
and then whole-frame sampling.
% ,
% a standard technique that runs upstream of the three \sys{} operators
% and is swept as an additional configuration parameter (\autoref{sec:additional-knobs}).
We also validate the one-dimensional tolerance sweep against the full combinatorial per-tile gap search.
% \subsubsection{Tracking Algorithm Selection}
% \label{sec:ablation-algorithm}
% To evaluate the impact of including tracking algorithm as a dimension of the Pareto search,
% we compare three configurations:
% (1)~only BYTETrack configurations,
% (2)~only SORT configurations,
% and (3)~the full Pareto frontier,
% which admits configurations from both algorithms.
% Depending on the dataset and the desired throughput-accuracy operating point,
% either algorithm may contribute Pareto-optimal points,
% so restricting the search to a single algorithm leaves gaps in the frontier.
% \ck{todo: add table or figure showing per-dataset Pareto frontiers under each condition, highlighting regimes where each algorithm contributes Pareto-optimal points and where the combined frontier strictly dominates.}

\subsubsection{Relevance Classifier Accuracy}
\label{sec:ablation-classifier}

We compare the deployed relevance classifier against the original RGB-only ShuffleNet baseline.
% on the \PtwozerofiveCompareClassifiersVideoset{} split.
Across \PtwozerofiveCompareClassifiersDatasetCount{} datasets at $T_r=\PtwozerofiveCompareClassifiersThreshold$,
adding the frame-difference image and position encoding improves F1 over the baseline on all \PtwozerofiveCompareClassifiersImprovedFOneDatasetCount{} datasets,
as shown in the \panellabel{Left} panel of \autoref{fig:classifier-and-packing}.
The augmented classifier yields an average F1 gain of \PtwozerofiveCompareClassifiersMacroAvgFOneDeltaPp{}\%.
The largest gain appears on \PtwozerofiveCompareClassifiersBestDatasetDisplay{},
where F1 rises from \PtwozerofiveCompareClassifiersBestFOneBaseline{} to \PtwozerofiveCompareClassifiersBestFOneModified{}
(gain \PtwozerofiveCompareClassifiersBestFOneDeltaPp{}\%).
This gain comes with \PtwozerofiveCompareClassifiersMacroAvgRuntimeSlowerPct{}\% higher classifier inference time on average.
\revpending{On three clips per dataset,
matched-frame timing on an idle Quadro RTX 6000 with local CPU/memory placement finds classification \revisionClassifierSpeedupMin{}--\revisionClassifierSpeedupMax{}$\times$ faster than detection at deployed batch sizes (16 versus one).
At equal batch size 16 it is \revisionClassifierEqualBatchSpeedupMin{}--\revisionClassifierEqualBatchSpeedupMax{}$\times$ faster;
including RGB copying gives \revisionClassifierWithCopySpeedupMin{}--\revisionClassifierWithCopySpeedupMax{}$\times$.
Classification achieves \revisionClassifierFpsMin{}--\revisionClassifierFpsMax{} FPS versus \revisionDetectorFpsMin{}--\revisionDetectorFpsMax{} detector FPS,
and occupies \revisionClassifierShareMinPct{}--\revisionClassifierShareMaxPct{}\% of extraction (median \revisionClassifierShareMedianPct{}\%) at the 5\% points.}
% We therefore use the augmented classifier in the end-to-end system.

\begin{figure}[t]
    \centering
    \includegraphics[width=\columnwidth]{figures/generated/p206_classifier_and_packing.pdf}
    \caption{\revpending{\panellabel{Left} Relevance-classifier comparison.}
    % across all
    % % \PtwozerofiveCompareClassifiersDatasetCount{}
    % datasets.
    Modified:
    \autoref{sec:relevance-classification}.
    % ShuffleNet augmented with the frame-difference image and tile-position encoding.
    Baseline:
    RGB-only ShuffleNet.
    % We report F1, recall, and precision against classifier runtime.
    % tile size \PtwozerofiveCompareClassifiersBestTileSize{},
    % and sample rate \PtwozerofiveCompareClassifiersSampleRate{}.
    \revpending{\panellabel{Right} Packing efficacy.}
    % ,
    % on the \packingEfficacyEvalSplit{}~split
    % across datasets,
    % at sample rate \packingEfficacySampleRate{},
    % restricted to the ``none'' tile-padding mode.
    % Packing efficacy is the fraction of total canvas tile-cells that the packer uses,
    % counting unoccupied tiles in the denominator.
    }
    \Description{}
    \label{fig:classifier-and-packing}
\end{figure}

% \subsubsection{pruning efficacy}
% \ck{show how our pruning algorithm is able to prune out at 0.5 sec runtime constraints}

\subsubsection{Packing Efficacy}
\label{sec:ablation-packing}
We quantify how tightly \sys{}'s packer fills each generated canvas by measuring the
fraction of total canvas tile-cells that the packer actually uses,
\[
\text{packing efficacy} \;=\; \frac{\text{occupied tiles}}{\text{occupied tiles} + \text{empty tiles}},
\]
where the denominator is
the number of canvases $\times$ canvas height $\times$ canvas width,
measured in tile units.
% This metric isolates the quality of the packing algorithm and is independent of the
% per-polyomino tile padding,
% which is intended to improve downstream detection rather than packing density.
% At sample rate \packingEfficacySampleRate{},
% and canvas scale \packingEfficacyCanvasScale{},
Without tile padding, the packer utilizes
\packingEfficacyMeanAcrossDatasets{}\% of the available canvas area on average across
\packingEfficacyDatasetCount{} datasets,
peaking at \packingEfficacyBestValue{}\% on \packingEfficacyBestDatasetDisplay{}.
The \panellabel{Right} panel of \autoref{fig:classifier-and-packing} shows the per-dataset packing
efficacy without tile-padding.

The two least densely packed datasets are B3D~2 and CalDoT~1,
which fall noticeably below the per-dataset average
(CalDoT~1 down to \packingEfficacyWorstValue{}\%).
We traced both to dataset-specific properties of the relevant polyominoes.
% rather than to limitations of the packer itself.
B3D~2 captures a busy unsignalized intersection in which vehicles routinely cluster
shoulder-to-shoulder.
% while waiting for
% % or yielding to
% crossing traffic.
Adjacent tiles with vehicles merge into a single connected polyomino,
so the per-canvas polyomino population is shifted toward fewer,
larger polyominoes than on the less-congested B3D scenes.
Specifically, aggregated across all validation videos,
B3D~2 produces \packingEfficacyBThreeDTwoPolyPerCanvas{} polyominoes per canvas
averaging \packingEfficacyBThreeDTwoPolyMeanTiles{} tiles each,
versus \packingEfficacyBThreeDThreePolyPerCanvas{} polyominoes per canvas
averaging \packingEfficacyBThreeDThreePolyMeanTiles{} tiles each on B3D~3.
% A handful of these merged blobs span the full canvas height,
% which blocks entire row bands of the canvas once placed.
CalDoT~1 exhibits a similar effect.
Aggregated across all validation videos,
CalDoT~1 produces \packingEfficacyCalDoTOnePolyPerCanvas{} polyominoes per canvas
averaging \packingEfficacyCalDoTOnePolyMeanTiles{} tiles each,
versus \packingEfficacyCalDoTTwoPolyPerCanvas{} polyominoes per canvas
averaging \packingEfficacyCalDoTTwoPolyMeanTiles{} tiles each on CalDoT~2.
% The CalDoT~1 polyominoes are also markedly taller and closer to square than on CalDoT~2,
% so they occupy both row and column slots of the fixed canvas
% rather than stacking cleanly along the bottom rows.
% Because \sys{}'s packer does not rotate polyominoes,
% since rotation would invalidate the
% tile-position encoding used by the relevance classifier
% (\autoref{sec:ablation-classifier}),
% these
% two effects translate directly into reduced packing density.

\revpending{Packing density measures canvas use;
its effect on detector accuracy is measured separately in \autoref{sec:packed-detection}.}

\subsubsection{Optimization Operators}
We measure the impact of each optimization technique by incrementally enabling them and comparing their throughput over the reference baseline under a strict 5\% maximum HOTA loss.
% Starting from the naive execution,
\sys{} first introduces the relevance classification and polyomino packing operators.
As shown by the Pareto front results,
this combined spatial pruning approach brings the throughput speedup over the reference execution up to \ablationOneSpeedupMax{}$\times$
(average \ablationOneSpeedupAvg{}$\times$)
while safely remaining within the tracking accuracy constraint.
Next, \sys{} adds the fine-grained polyomino pruning operator, which significantly accelerates tracking in simple, low-activity video regions.
The inclusion of polyomino pruning further increases the maximum speedup to \ablationTwoSpeedupMax{}$\times$
(average \ablationTwoSpeedupAvg{}$\times$)
over the baseline execution.
Finally, we evaluate enabling the whole-frame sampling knob $s > 1$ (\autoref{sec:additional-knobs}).
% for high-throughput operating points where the accuracy budget permits skipping whole frames upstream of the tracking pipeline.
Together,
the complete execution engine achieves a maximum speedup of \ablationThreeSpeedupMax{}$\times$
(average \ablationThreeSpeedupAvg{}$\times$)
while still strictly adhering to the 5\% tracking accuracy loss constraint.

% \subsubsection{why canvas size has to be the same as the video frame size?}
% \subsubsection{Canvas Size}
% \label{sec:ablation-canvas-size}
% \sys{} fixes the canvas pixel dimensions to match those of the source video frame
% (canvas scale \packingEfficacyCanvasScale{}) for three reasons.
% First,
% both the relevance classifier and the object detector are invoked on canvas content
% rather than on raw frames.
% Each was trained on inputs drawn at the source-frame scale,
% so a canvas at the same scale preserves the per-tile pixel resolution they expect
% and avoids any train/test distribution shift caused by re-sampling.
% Second,
% polyominoes can be arbitrarily large within the always-relevant region of a frame.
% Empirically,
% the bounding box of the largest polyominoes we observed fills nearly the full source
% frame;
% for example,
% $12 \times 15$ tiles inside a $12 \times 18$-tile canvas on B3D~2.
% A canvas strictly smaller than the source frame would force such polyominoes to be
% split,
% breaking the tile-position encoding used by the classifier
% (\autoref{sec:ablation-classifier})
% and complicating the inverse mapping from detected boxes back to source-frame
% coordinates.
% Third,
% a canvas strictly larger than the source frame would waste detector compute on tiles
% that can never be filled,
% since the foreground area per frame is upper-bounded by the source-frame area itself.
% Setting the canvas equal to the source frame size therefore upper-bounds the packing
% problem at exactly the resolution required for downstream tracking accuracy,
% with no excess detector work.

\subsubsection{Maximum tile sampling gaps}
\label{sec:ablation-constraints}

We evaluate the sweep over the mistrack-rate tolerance~$\bar{M}$ from \autoref{sec:learn-constraints}.
Each tolerance yields a different $\bar{G}^{\bar{M}}$.
We show that the sweep approximates the Pareto frontier of pruning ratio versus HOTA against obtained by exhaustively enumerating all possible $\bar{G}$.
% exhaustive per-tile gap selection.
To keep the exhaustive baseline feasible,
we reduce the tile grid to $3 \times 3$
and the candidate gap set to $\{1, 2, 4\}$,
yielding $3^9 \approx 20{,}000$ per-tile gap assignments.
We do not report throughput directly on this sweep,
because re-running full tracking execution on every assignment is prohibitively expensive;
instead,
we use the pruning ratio (fraction of tiles skipped)
as a throughput proxy.
We use BYTETrack as the tracker.

To isolate the effect of the maximum tile sampling gaps
from classifier and detector variance,
the experimental pipeline constructs polyominoes directly from the reference-pipeline bounding boxes rather than classifier output,
retains only detections whose centers fall inside a retained polyomino,
and feeds the retained detections to BYTETrack.
We report HOTA against the reference tracking output
alongside the pruning ratio.

\begin{figure}[t]
    \centering
    \includegraphics[width=\columnwidth]{figures/generated/pruning_vs_hota.png}
    \caption{HOTA vs.\ pruning ratio per dataset.
    Gray: exhaustive per-tile sampling assignments
    over a $3\times 3$ grid with gaps $\{1, 2, 4\}$.
    Black: Pareto frontier of the exhaustive sweep.
    Blue diamonds: our heuristic $\bar{G}^{\bar{M}}$ swept over $\bar{M}$ with BYTETrack.}
    \Description{}
    \label{fig:sampling-constraints-ablation}
    % \vspace{-0.25in}
\end{figure}

\autoref{fig:sampling-constraints-ablation} shows the resulting trade-off per dataset.
The heuristic points (blue diamonds),
obtained by sweeping $\bar{M}$,
lie on or near the Pareto frontier of the exhaustive per-tile sweep across datasets,
indicating that the one-dimensional tolerance sweep loses little
relative to the full combinatorial search.
Quantitatively,
at each heuristic pruning level we compare against the exhaustive
Pareto point with the highest pruning ratio not exceeding that level;
HOTA loss versus that anchor is at most \MistrackHeuristicMaxHotaLossPercent\%{}
(mean \MistrackHeuristicAvgHotaLossPercent\%).
\input{content/revision/additional_evaluation}
